
\section*{Part II. Applications of optics \hfill 7 / 15}
\subsection*{1. A noisy optical device is an open quantum system}
For a system \(S\) interacting with an environment \(E\), evolution
is unitary on the \emph{joint} space but generally not on \(S\):
\[
\boxed{\rho'_S=\mathcal E(\rho_S)
=\operatorname{Tr}_E[U(\rho_S\otimes\rho_E)U^\dagger].}
\]
The same channel can be expressed with Kraus operators
\(\mathcal E(\rho)=\sum_kK_k\rho K_k^\dagger\),
where \(\sum_kK_k^\dagger K_k=I\) for trace preservation.
The measurement outcome \(k\) is ignored here.
For pure attenuation, mix the signal \(a\) with an
environment vacuum mode \(e\) on a beam splitter:
\[
\hat a'=\sqrt\eta\,\hat a+\sqrt{1-\eta}\,\hat e,
\quad
\ket{1,0}\to\sqrt\eta\ket{1,0}
+\sqrt{1-\eta}\ket{0,1}
\]
(up to the selected phase convention).
Tracing out the environment gives
\[
\boxed{\mathcal E_\eta(\ket1\bra1)
=\eta\ket1\bra1+(1-\eta)\ket0\bra0.}
\]
A single-photon superposition loses coherence as well:
\(\mathcal E_\eta(\ket0\bra1)=\sqrt\eta\ket0\bra1\).
One must not just replace \(\hat a\to\sqrt\eta\hat a\):
the vacuum noise operator preserves commutators.

\subsection*{2. Continuous time: Lindblad equation}
A Markovian reservoir produces
\[
\boxed{\dot\rho=-\frac{i}{\hbar}[H,\rho]
+\sum_j\gamma_j\mathcal D[L_j]\rho},\quad
\mathcal D[L]\rho=L\rho L^\dagger
-\tfrac12\{L^\dagger L,\rho\}.
\]
For cavity photon loss take \(L=\hat a\). The average photon
number decays as \(\langle n(t)\rangle=\langle n(0)\rangle e^{-\gamma t}\)
for zero-temperature damping. A thermal bath requires absorption
as well as emission channels.

\subsection*{3. Langevin equation: damping must accompany noise}
In the Heisenberg picture a cavity with frequency \(\omega_c\)
and damping rate \(\kappa\) follows a form such as
\[
\dot{\hat a}=-(i\omega_c+\kappa/2)\hat a
+\sqrt\kappa\,\hat a_{\rm in}(t)
\]
with sign convention set by the bath coupling.
The input-noise operator is necessary to retain the
canonical commutator. In the Gardiner--Collett derivation,
a continuum of reservoir modes, a near-resonant/Markov
approximation and bath initial conditions lead to both the
damping and fluctuation terms.

\begin{methodbox}
\textbf{Method:} choose \(S,E\); specify environment state and
unitary interaction; apply it to the joint input; trace over
\(E\). For time evolution, identify jump operators and rates.
Always check \(\operatorname{Tr}\rho'=1\).
\end{methodbox}
\src{W5L1 pp.~7--31; W5L2 pp.~2--25; W2L2 pp.~15--17.}
